\documentclass[11pt]{article}

% Prefer packages supported by arXiv's TeX Live and LaTeXML environments.
\usepackage[T1]{fontenc}
\usepackage{lmodern}
\usepackage{microtype}
\usepackage{amsmath,amssymb}
\usepackage{graphicx}
\usepackage{xurl}
\usepackage[hidelinks]{hyperref}
\hypersetup{pdftitle={voiage: Value of Information Analysis with a Rust Core and Polyglot Bindings},
  pdfauthor={Dylan Mordaunt},
  pdfsubject={Open-source research software for value of information analysis},
  pdfkeywords={value of information, decision analysis, health economics, uncertainty quantification, Rust, Python}
}

\title{\texttt{voiage}: Value of Information Analysis with a Rust Core and Polyglot Bindings}
\author{Dylan Mordaunt\\\small Independent Researcher, Australia}
\date{}

\begin{document}
\maketitle

\begin{abstract}
Value of Information (VOI) analysis estimates how much a decision could
improve if uncertainty were reduced by collecting additional information.
\texttt{voiage} is open-source research software for computing and
communicating VOI analyses from probabilistic decision models. Version 1.0.0
combines a Python user interface with authoritative Rust numerical kernels and
shared C-ABI bindings for R and Julia. The software implements expected value
of perfect information, expected value of partial perfect information,
expected value of sample information, expected net benefit of sampling,
cost-effectiveness acceptability frontiers, dominance, and related diagnostics.
Language-neutral fixtures, explicit schemas, and layered assurance gates make
the numerical contract reproducible across supported bindings.
\end{abstract}

\section{Summary}

Value of Information (VOI) analysis asks which uncertainty matters to a
decision, what evidence could reduce it, and whether the expected improvement
justifies the cost of learning. In health technology assessment this may mean
choosing a trial, targeting a subgroup, validating a prediction model, or
improving uptake. The formulation can also be used for environmental
monitoring, engineering reliability, manufacturing quality, energy planning,
and investment decisions when uncertain consequences are mapped to a common
decision-value scale.

The difficult part is usually not the textbook expected value of perfect
information (EVPI) equation. Analysts must decide what counts as a strategy,
whose outcomes matter, which population and time horizon to scale to, whether a
model form is credible, whether evidence transports to the target setting, and
whether a recommendation will be implemented. Different choices can change
the preferred action even when the same simulation output is used. The
\texttt{voiage} package records these choices in the analysis inputs,
diagnostics, and returned result rather than leaving them as unrecorded
preprocessing.

The implemented elements described in this paper are:
\begin{itemize}
  \item a stable analysis path from probabilistic net benefits to
  EVPI, expected value of partial perfect information (EVPPI), expected value
  of sample information (EVSI), expected net benefit of sampling (ENBS),
  dominance, the cost-effectiveness acceptability frontier (CEAF),
  heterogeneity, and structural VOI\@;
  \item a labelled decision representation that records sample, strategy,
  parameter, subgroup, perspective, population, and study-design semantics;
  \item a maturity-labelled frontier that represents value questions
  about equity, perspective, implementation, validation, transportability,
  data quality, computation, and staged decisions; and
  \item diagnostic result envelopes that retain method, seed, warning,
  scaling, maturity, and provenance information that supports interpretation
  of a value estimate.
\end{itemize}

Version 1.0.0 routes selected stable aggregation kernels through
binding-independent Rust components. Python retains model orchestration,
validation for broader method paths, labelled data, plots, and
reports~\cite{voiage2026}. The R and Julia bindings currently expose direct
native EVPI only; wider application programming interface (API) parity is not
claimed. Definitions and abbreviations used throughout the paper are collected in
Sections~\ref{sec:glossary} and~\ref{sec:abbreviations}.

\section{Statement of need}

VOI analysis is central to deciding whether additional research is worthwhile,
but practical implementations are often fragmented across domain-specific
packages, languages, and model formats. This fragmentation complicates
comparison, reproduction, edge-case validation, and movement from a research
prototype to an operational decision workflow. Python users in particular need
a maintained, open implementation that supports common VOI methods while
remaining compatible with the scientific Python ecosystem.

\texttt{voiage} addresses this need with a common decision-model contract,
validated schemas, deterministic fixtures, explicit diagnostics, reproducible
seeds, machine-readable result envelopes, and a command-line interface. It is
designed for researchers and analysts evaluating proposed data collection,
comparing decisions under uncertainty, or integrating VOI calculations into
larger evidence and policy pipelines.

\section{State of the field}

VOI is established in statistical decision analysis and health-economic
evaluation~\cite{claxton1999irrelevance,briggs2006decision}. R tooling
includes the Bayesian cost-effectiveness analysis (BCEA)
ecosystem~\cite{baio2017bcea}, while methodological work has made EVPPI and EVSI more
practical using regression estimators~\cite{strong2014evppi,strong2015evsi}.
Methods also exist for structural uncertainty, heterogeneity, distributional
cost-effectiveness, implementation, and adaptive trial
design~\cite{jackson2011structural,espinoza2014heterogeneity,asaria2016distributional,fenwick2008implementation,flight2022sequential}.
These are complementary advances. Reimplementing their names does not by
itself constitute a methodological contribution, and \texttt{voiage} does not
claim statistical superiority over specialist implementations.

These concerns can be evaluated in separate analyses even when they act on the
same decision. Research may have little
value if implementation is weak; subgroup optimization can conflict with an
equity-weighted objective; evidence can be precise in a source population but
poorly transportable; and waiting can preserve option value when commitment is
costly to reverse. Treating each as unrelated post-processing makes it hard to
compare their assumptions or combine them in an evidence plan.

The \texttt{voiage} implementation uses a shared, maturity-labelled problem
architecture for these value questions. Stable
classical estimands and fixture-backed frontier methods use common strategy,
value, diagnostic, and reporting semantics. The shared representation is
intended to support questions about whether the limiting uncertainty is a
parameter, study, structure, subgroup, perspective, implementation process,
data source, or decision stage without changing the basic representation of
the decision. This is an architectural design choice, not evidence of
methodological superiority. Several frontier formulations still require
empirical validation, methodological comparison, and cross-language parity
before stable promotion.

\section{Software design and assurance}

The stable runtime is organized around validated domain objects, canonical
schemas, numerical kernels, diagnostics, provenance, and adapters. The Rust
workspace contains the authoritative numerical implementation and a narrow,
versioned C application binary interface with explicit ownership and error
transport. The Python package
uses PyO3 and Maturin for its native extension and retains the facade and
workflow responsibilities needed by users. R and Julia consume the same
interface;
their packages do not reimplement VOI policy. Optional dependencies such as
JAX, plotting libraries, and web integrations are isolated from the base core.

Correctness is supported by language-neutral fixtures, unit and integration
tests, property-based and fuzz tests, metamorphic checks, mutation testing,
coverage gates, Rust diagnostics, ABI checks, benchmark regression budgets,
and clean-install tests. The release process produces platform wheels, a
source archive, checksums, a software bill of materials, provenance evidence,
and a signed GitHub release. Documentation is maintained in Astro/Starlight.
The v1.0.0 release is available from GitHub and the Python Package Index, while
external registry and
curation outcomes remain tracked separately.

The manuscript follows software-citation guidance~\cite{smith2016software}.
Its canonical source is authored LaTeX. The repository builds a deterministic
source-only archive, audits the review PDF, compiles against the supported
arXiv TeX Live generations, and independently converts the source to a
semantic web document. These checks establish repository readiness, not submission or
acceptance by arXiv or a journal.

\section{Availability, reproducibility, and intended use}

\texttt{voiage} is intended to support research in health technology
assessment, clinical-trial design, public policy, environmental
decision-making, and related areas of uncertainty quantification. The
repository contains deterministic analytical and integration fixtures,
tutorials, cross-language contracts, and representative workflows for these
use cases.

Version 1.0.0 is available under the Apache License 2.0 from
\url{https://github.com/edithatogo/voiage/releases/tag/v1.0.0} and
\url{https://pypi.org/project/voiage/1.0.0/}.
The repository records citation metadata, supported Python versions, the Rust
minimum supported version, locked dependencies, release checksums, and
machine-readable provenance. The documentation site is built from the same
repository using Astro/Starlight. Examples and normative fixtures use synthetic
or explicitly governed inputs so that the package can be evaluated without
access to confidential decision models.

The project applies the principle that research software should be findable,
accessible, interoperable, and reusable where technically
possible~\cite{wilkinson2016fair}. In practice, this means stable identifiers for
schemas, portable fixtures, explicit licenses, documented installation paths,
and result envelopes that can be inspected outside the originating language.
It does not mean that every experimental method or external registry has the
same maturity.

Potential applications include comparing research designs, prioritizing
parameters for further study, assessing decision uncertainty across
subgroups, and embedding VOI results in larger evidence pipelines. These are
intended uses, not evidence of adoption or policy impact. Publicly attributable
case studies can be added in a later manuscript revision when such evidence is
available.

\section{Limitations}

The breadth of the framework does not make value functions interchangeable
across domains. Users remain responsible for model structure, causal
assumptions, input provenance, outcome valuation, thresholds, distributional
choices, and interpretation. A monetary net-benefit model, an equity-weighted
social objective, and an engineering loss function can share computational
structure while representing different ethical and empirical claims.

Regression-based EVPPI and EVSI estimates inherit modelling choices and Monte
Carlo error. Convergence, effective sample size, predictive adequacy,
extrapolation, and sensitivity to method settings provide useful context for a
scalar result. Population VOI additionally depends on
eligible population, time horizon, discounting, implementation, and research
cost assumptions; \texttt{voiage} validates representations and numerical
preconditions but cannot determine whether those assumptions are substantively
appropriate.

Newer methods are individually marked as developing or experimental. Their
repeatable examples show that the stated calculations and software behaviour
have been implemented. They do not establish statistical validity, accuracy in
real applications, superiority to other methods, suitability in another
setting, or usefulness for policy. In particular,
perspective, data-quality, transportability, dynamic-option, and several
evidence-system methods need applied comparisons and domain-specific validation.
The Python package is broader than the current R and Julia packages. Both
provide native EVPI and ENBS in version 2.2.0, but neither exposes Python's
full decision-record interface. R's optional EVPPI and EVSI paths require Python.

Automated checks detect some implementation errors but do not replace
independent statistical validation, expert review, sensitivity analysis, or
replication with relevant domain data. Optional accelerators and distributed
or hardware research are not evidence of faster production analysis.

\section{AI usage disclosure}

Generative AI tools assisted with repository analysis, security and release
workflow review, documentation drafting, and manuscript preparation. The human
author remains responsible for the software design, claims, references, author
list, and final text. All AI-assisted changes are subject to human review and
repository tests; no AI tool is an author or a substitute for scientific and
editorial judgment.

\section{Acknowledgements}

The author acknowledges the open-source communities whose statistical,
scientific-computing, and programming tools support this work. Funding,
institutional support, conflicts of interest, and any additional contributors
remain for the human author to declare before submission.


\bibliographystyle{plain}
\bibliography{references}

\end{document}
